\chapter{Plancherel 定理}
\label{chap:plancherel-theorem}

\(L^1\) Fourier 变换给出有界连续函数，却不保持 \(L^1\) 范数，也未必产生
可积函数。能量计算使用平方积分范数；在这一范数下，Fourier 变换成为酉算子。
证明不能直接把两个无穷积分随意换序。我们先在
\(C_c^\infty(\R^d)\) 上利用反演定理建立内积恒等式，再借助稠密性把变换延拓
到整个 \(L^2\)。

\section{一个稠密的试验类}

由\chapterref{chap:lp-approximation}，\(C_c^\infty(\R^d)\) 在
\(L^2(\R^d)\) 中稠密。若 \(\varphi\in C_c^\infty\)，则
\(\widehat\varphi\in L^1\cap L^2\)。为看出无穷远衰减，任选整数
\(m>d/2\)，对算子 \((I-\Delta)^m\) 分部积分得
\[
 (1+4\pi^2|\xi|^2)^m\widehat\varphi(\xi)
 =\widehat{(I-\Delta)^m\varphi}(\xi).
\]
右侧由 \(L^1\) 范数一致控制，所以
\(|\widehat\varphi(\xi)|\le C(1+|\xi|)^{-2m}\)，足以同时属于
\(L^1\) 和 \(L^2\)。

\begin{proposition}[Parseval 公式]\label{prop:fourier-parseval-test}
若 \(\varphi,\psi\in C_c^\infty(\R^d)\)，则
\[
 \int_{\R^d}\widehat\varphi(\xi)\overline{\widehat\psi(\xi)}\,d\xi
 =\int_{\R^d}\varphi(x)\overline{\psi(x)}\,dx.
\]
\end{proposition}
\begin{proof}
\(\psi,\widehat\psi\) 都可积，反演定理给
\[
 \overline{\psi(x)}
 =\int_{\R^d}\overline{\widehat\psi(\xi)}e^{-2\pi ix\cdot\xi}\,d\xi.
\]
乘以 \(\varphi(x)\) 后的二重积分绝对可积，因为
\(\|\varphi\|_1\|\widehat\psi\|_1<\infty\)。Fubini 定理于是给
\[
 \int\varphi(x)\overline{\psi(x)}dx
 =\int\overline{\widehat\psi(\xi)}
      \left(\int\varphi(x)e^{-2\pi ix\cdot\xi}dx\right)d\xi.
\]
\end{proof}

取 \(\psi=\varphi\)，得到 \(\|\widehat\varphi\|_2=\|\varphi\|_2\)。

\section{酉延拓}

\begin{theorem}[Plancherel 定理]\label{thm:plancherel}
存在唯一酉算子
\[
 \mathcal F:L^2(\R^d)\longrightarrow L^2(\R^d)
\]
在 \(L^1\cap L^2\) 上与积分定义的 Fourier 变换一致。其逆算子为
\[
 \mathcal F^{-1}g(x)=\mathcal Fg(-x)
\]
在稠密试验类上的连续延拓，并且
\[
 \langle\mathcal Ff,\mathcal Fg\rangle_{L^2}
 =\langle f,g\rangle_{L^2},\qquad
 \|\mathcal Ff\|_2=\|f\|_2.
\]
\end{theorem}
\begin{proof}
若 \(f\in L^2\)，取 \(\varphi_n\in C_c^\infty\) 使
\(\varphi_n\to f\) 于 \(L^2\)。Parseval 公式给
\(\|\widehat\varphi_n-\widehat\varphi_m\|_2
=\|\varphi_n-\varphi_m\|_2\)，故变换后的序列在 \(L^2\) 中收敛。定义
\(\mathcal Ff=\lim\widehat\varphi_n\)。相同估计说明定义与逼近列无关，且
得到线性和等距性。

在使用两次变换之前，须核对这个延拓与积分定义的一致性。若
\(f\in L^1\cap L^2\)，先用 \(f\mathbf1_{B(0,n)}\) 截断，再选足够小的
磨光参数，可得 \(\varphi_n\in C_c^\infty\) 同时在 \(L^1,L^2\) 中趋于 \(f\)。
于是 \(\widehat{\varphi_n}\to\widehat f\) 一致，而按延拓的定义，
\(\widehat{\varphi_n}\to\mathcal Ff\) 于 \(L^2\)。后一收敛有几乎处处收敛
子列，故 \(\mathcal Ff=\widehat f\) 几乎处处。这里截断与磨光在两种范数
中的误差都可分别控制在 \(1/n\) 内，所以不需要为两种范数选择不同的最终序列。

在试验类上，两次变换满足
\(\mathcal F^2\varphi(x)=\varphi(-x)\)：因为
\(\widehat\varphi\in L^1\cap L^2\)，第二次延拓变换也可以用积分计算，
此时再应用反演定理。反射算子
\(R\varphi(x)=\varphi(-x)\) 在 \(L^2\) 上酉，故稠密延拓给
\(\mathcal F^2=R\)。于是 \(\mathcal F\) 满射，并且
\(\mathcal F^{-1}=R\mathcal F=\mathcal FR\)。极化公式再给内积保持。
\end{proof}

证明中核对两种定义的步骤不可省略。\(\widehat\varphi\) 通常没有紧支撑，
因此不能仅凭 \(\varphi\in C_c^\infty\)，就把第二次延拓变换当作积分变换。
一致性一经建立，今后在 \(L^1\cap L^2\) 上便可按计算需要使用任一定义。

\section{\texorpdfstring{\(L^2\)}{L2} 反演与截断}

对一般 \(f\in L^2\)，表达式
\(\int f(x)e^{-2\pi ix\cdot\xi}dx\) 未必绝对收敛，不能把它当作逐点定义。
Plancherel 变换是 \(L^2\) 等价类的范数极限。若
\[
 f_R=f\,\mathbf1_{B(0,R)},
\]
则 \(f_R\in L^1\cap L^2\)，且 \(f_R\to f\) 于 \(L^2\)，所以
\[
 \widehat{f_R}\longrightarrow\mathcal Ff\qquad\text{于 }L^2.
\]
不同截断方式只要在 \(L^2\) 中逼近 \(f\)，都会得到同一个等价类；逐点极限
则可能不存在。

同样，若 \(g=\mathcal Ff\)，对频率截断
\(g_R=g\mathbf1_{B(0,R)}\) 作逆变换，得到带限函数
\[
 f_R^\sharp(x)=\int_{B(0,R)}g(\xi)e^{2\pi ix\cdot\xi}\,d\xi,
\]
并有 \(f_R^\sharp\to f\) 于 \(L^2\)。这是 \(L^2\) 反演的可计算形式。

\section{卷积与乘积的平方积分版本}

若 \(f\in L^2,g\in L^1\)，先证明所需的 Young 估计。\(g=0\) 时显然。
否则对概率测度 \(|g(y)|dy/\|g\|_1\) 使用 Cauchy--Schwarz，得到
\[
 |f*g(x)|^2\le\|g\|_1\int|f(x-y)|^2|g(y)|\,dy.
\]
Tonelli 定理说明右侧对 \(x\) 的积分为 \(\|g\|_1^2\|f\|_2^2<\infty\)。
因此内层积分几乎处处有限，再用同一概率测度的 Cauchy--Schwarz 可知
卷积的绝对积分几乎处处有限。积分上述估计便得
\(\|f*g\|_2\le\|f\|_2\|g\|_1\)。用 \(L^1\cap L^2\) 函数逼近 \(f\)，
从 \(L^1\) 卷积定理和 Plancherel 连续性得到
\begin{equation}\label{eq:plancherel-convolution}
 \mathcal F(f*g)=(\mathcal Ff)\widehat g\qquad\text{于 }L^2.
\end{equation}
乘以有界频率函数 \(\widehat g\) 因而对应一个有界平移不变算子。

若 \(f,g\in L^2\)，则 \(fg\in L^1\)，而两个 \(L^2\) 函数的卷积可由
绝对收敛积分定义：
\[
 ((\mathcal Ff)*(\mathcal Fg))(\xi)
 =\int_{\R^d}(\mathcal Ff)(\xi-\eta)(\mathcal Fg)(\eta)\,d\eta.
\]
Cauchy--Schwarz 不等式给这个积分以 \(\|f\|_2\|g\|_2\) 为共同上界。
具体取 \(f_n,g_n\in C_c^\infty\)，分别在 \(L^2\) 中逼近 \(f,g\)。
它们及其变换均在 \(L^1\) 中，故可使用上一章的乘积公式。又有
\[
 \|f_ng_n-fg\|_1
 \le\|f_n-f\|_2\|g_n\|_2+\|f\|_2\|g_n-g\|_2\longrightarrow0.
\]
在频率端用相同的 Cauchy--Schwarz 估计，卷积的差一致地趋零。
因此两边都能一致地传到极限，得到
\[
 \widehat{fg}=(\mathcal Ff)*(\mathcal Fg).
\]
左侧属于 \(C_0(\R^d)\)，所以等式也证明右侧连续且在无穷远消失。
这里的卷积属于 \(L^\infty\)，未必属于
\(L^1\)。

\section{微分算子的频率定义}

对 \(f\in L^2(\R^d)\)，若
\(\xi_j\mathcal Ff(\xi)\in L^2\)，定义
\[
 D_jf=\mathcal F^{-1}\bigl(2\pi i\xi_j\mathcal Ff\bigr).
\]
其定义域为
\[
 \mathcal D(D_j)=
 \{f\in L^2:\xi_j\mathcal Ff\in L^2\}.
\]
若 \(f\in C_c^\infty\)，它与经典偏导数相同。因此定义域包含
\(C_c^\infty\)，在 \(L^2\) 中稠密。
称这样的算子闭，是指其图在 \(L^2\times L^2\) 中闭。
用序列检验：若
\(f_n\to f\) 且 \(D_jf_n\to g\) 于 \(L^2\)，对等式
\(\mathcal F(D_jf_n)=2\pi i\xi_j\mathcal Ff_n\)，选一个使两侧涉及的
两个 \(L^2\) 收敛列同时几乎处处收敛的子列，再逐点取极限，可知
\(f\in\mathcal D(D_j)\) 且 \(D_jf=g\)，所以算子闭。

这并不是定义在整个 Banach 空间上的闭图定理的情形。
例如取单位体积的频率长方体，在第 \(j\) 个坐标上为 \([n,n+1]\)，
其余坐标为 \([0,1]\)，以其指标函数作为 \(\widehat f_n\)。则
\(f_n\in\mathcal D(D_j)\)、\(\|f_n\|_2=1\)，却有
\(\|D_jf_n\|_2\ge2\pi n\)。所以 \(D_j\) 不满足关于 \(L^2\) 范数的有界估计。

为处理下一节的一维问题，先建立此时频率导数与函数导数的联系。
对 \(f,g\in L^2(\R)\)，称 \(g\) 是 \(f\) 的弱导数，如果
\[
 \int_\R f\varphi'=-\int_\R g\varphi
 \qquad(\varphi\in C_c^\infty(\R)).
\]
这是周期情形的相同积分条件，只是测试函数现在具有紧支撑。

\begin{proposition}\label{prop:fourier-derivative-ac-line}
对 \(f,g\in L^2(\R)\)，下列条件等价：
\(f\in\mathcal D(D_1)\) 且 \(D_1f=g\)；
\(g\) 是 \(f\) 的弱导数；
\(f\) 有一个局部绝对连续代表 \(F\)，满足 \(F'=g\) 几乎处处。
\end{proposition}
\begin{proof}
若 \(D_1f=g\)，对测试函数 \(\varphi\) 用 Plancherel 得
\[
 \langle g,\varphi\rangle
 =\langle2\pi i\xi\widehat f,\widehat\varphi\rangle
 =-\langle\widehat f,2\pi i\xi\widehat\varphi\rangle
 =-\langle f,\varphi'\rangle.
\]
测试函数及其共轭属于同一类，故这正是弱导数条件。

若 \(g\) 是弱导数，令 \(G(x)=\int_0^xg(t)dt\)，并置 \(h=f-G\)。
\(G\) 局部绝对连续，分部积分给 \(\int h\varphi'=0\)。
对任意积分为零的 \(\psi\in C_c^\infty\)，其原函数
\(\varphi(x)=\int_{-\infty}^x\psi(t)dt\) 仍在 \(C_c^\infty\)，所以
\(\int h\psi=0\)。固定 \(\eta\in C_c^\infty\) 使 \(\int\eta=1\)，
记 \(c=\int h\eta\)。将任意测试函数减去 \((\int\psi)\eta\)，得到
\(\int(h-c)\psi=0\)。在每个有界开区间上，\(h-c\in L^2\)，而
\(C_c^\infty\) 在该区间的 \(L^2\) 中稠密（\exerciseref{density-open-set}），
故 \(h=c\) 几乎处处。
因此 \(F=G+c\) 是所需代表。

最后设有这样的代表 \(F\)。取固定光滑截断 \(\chi\)，在 \([-1,1]\)
上为一、支撑于 \([-2,2]\)、取值于 \([0,1]\)，令 \(\chi_R(x)=\chi(x/R)\)。
则 \(u_R=\chi_RF\) 绝对连续且紧支撑，
\[
 u_R\to f,\qquad u_R'=\chi_Rg+\chi_R'F\to g\quad\text{于 }L^2.
\]
第一项由控制收敛，第二式中的截断导数项由
\(\|\chi_R'F\|_2\le\|\chi'\|_\infty\|f\|_2/R\) 控制。
取非负、积分为一的紧支撑光滑磨光核 \(\rho_\varepsilon\)。
对固定 \(R\)，积分号下求导再对绝对连续函数分部积分，给
\((u_R*\rho_\varepsilon)'=u_R'*\rho_\varepsilon\)。
边界项因紧支撑消失；涉及的函数在有限区间上可积，故换序合法。
磨光逼近定理使函数与导数同时在 \(L^2\) 中趋于 \(u_R,u_R'\)。
依次取 \(R\to\infty\) 和足够小的 \(\varepsilon\)，得到
\(v_n\in C_c^\infty\)，使 \(v_n\to f\)、\(v_n'\to g\) 于 \(L^2\)。
已知 \(D_1v_n=v_n'\)，由 \(D_1\) 的闭性得到 \(D_1f=g\)。
\end{proof}

\chapterref{chap:distributions-weak-derivatives}将系统讨论开集上的分布导数
和弱导数；本章所需的一维等价关系已经由上述论证建立。

\section{一个不确定性估计}

若 \(f\in L^2(\R)\)、\(xf\in L^2\)，且 \(f\) 有平方可积弱导数，则
\[
 \|f\|_2^2=-2\operatorname{Re}\int_\R x f'(x)\overline{f(x)}\,dx
\]
先在光滑紧支撑类中由分部积分成立。下面核实推广所需的共同逼近。
使用上一命题的局部绝对连续代表和截断 \(u_R=\chi_Rf\)，除了函数及导数
收敛，还有 \(xu_R\to xf\) 于 \(L^2\)。对固定 \(R\)，磨光满足
\[
 x(u_R*\rho_\varepsilon)
 =(xu_R)*\rho_\varepsilon+u_R*(x\rho_\varepsilon).
\]
第一项在 \(L^2\) 中趋于 \(xu_R\)，第二项的范数不超过
\(\|u_R\|_2\|x\rho_\varepsilon\|_1=O(\varepsilon)\|u_R\|_2\)，
其中常数只由固定磨光核决定。于是可选 \(v_n\in C_c^\infty\)，使
\(v_n\to f\)、\(xv_n\to xf\)、\(v_n'\to f'\) 同时在 \(L^2\) 中成立。
积分乘积由 Cauchy--Schwarz 连续地传递到极限，所以上述恒等式对题设函数成立。
再用 Cauchy--Schwarz 给
\[
 \|xf\|_2\,\|f'\|_2\ge\frac12\|f\|_2^2.
\]
结合 \(\mathcal F(f')=2\pi i\xi\mathcal Ff\)，得到
\begin{equation}\label{eq:heisenberg-simple}
 \|xf\|_2\,\|\xi\mathcal Ff\|_2
 \ge\frac1{4\pi}\|f\|_2^2.
\end{equation}
零函数显然取等号。若 \(f\ne0\)，实部估计和 Cauchy--Schwarz 同时取等号
要求 \(f'=-cxf\) 几乎处处，其中 \(c>0\)。对局部绝对连续代表乘以
\(e^{cx^2/2}\)，其导数几乎处处为零，因而在每个紧区间上为常数。
故非零等号函数恰为 \(f(x)=Ce^{-a x^2}\)，其中 \(a>0\)、\(C\ne0\)；
直接代入也验证这些函数取等号。这里以原点和零频为基准。
若先把 \(x\) 和频率分别减去各自的
中心，再写相应的方差不等式，则平移、调制和伸缩后的 Gaussian 函数给出
全部等号情形。

\section*{习题}
\addcontentsline{toc}{section}{习题}
\subsection*{A组}
\begin{enumerate}[label=\textbf{\thechapter.\arabic*.}]
\exerciseitem{plancherel-test-decay}证明 \(C_c^\infty\) 函数的 Fourier 变换快于任意指定多项式阶衰减。
\exerciseitem{plancherel-inner-product}从反演定理严谨推导试验函数上的 Parseval 公式，并标出 Fubini 的控制量。
\exerciseitem{plancherel-extension-independent}证明等距延拓的定义与所选 \(C_c^\infty\) 逼近列无关。
\exerciseitem{plancherel-square-reflection}在试验函数上证明 \(\mathcal F^2f(x)=f(-x)\)，再推广到 \(L^2\)。
\end{enumerate}
\subsection*{B组}
\begin{enumerate}[label=\textbf{\thechapter.\arabic*.},resume]
\exerciseitem{plancherel-lone-agreement}证明 Plancherel 延拓在 \(L^1\cap L^2\) 上与积分定义一致。
\exerciseitem{plancherel-frequency-truncation}证明频率截断反演 \(f_R^\sharp\to f\) 于 \(L^2\)，并求误差范数。
\exerciseitem{plancherel-convolution-ltwo}证明公式\eqref{eq:plancherel-convolution}及其范数估计。
\exerciseitem{plancherel-derivative-closed}证明频率定义的 \(D_j\) 是闭算子，且 \(C_c^\infty\) 包含于其定义域。
\exerciseitem{plancherel-translation-generator}设 \(f\in L^2(\R^d)\)。证明
\((\tau_{he_j}f-f)/h\) 在 \(h\to0\) 时有 \(L^2\) 极限，当且仅当
\(\xi_j\mathcal Ff\in L^2\)；此时极限为 \(-D_jf\)。
\end{enumerate}
\subsection*{C组}
\begin{enumerate}[label=\textbf{\thechapter.\arabic*.},resume]
\exerciseitem{plancherel-heisenberg}补全不确定性估计\eqref{eq:heisenberg-simple}的逼近论证，并求等号条件。
\exerciseitem{plancherel-poisson-solver}设 \(f\in L^2(\R^d)\) 且 \(|\xi|^{-2}\widehat f(\xi)\in L^2\)。用 Fourier 变换构造 \(L^2\) 解 \(u\) 满足 \(-\Delta u=f\)，并说明零频附近的条件。这里 \(\Delta=\sum_jD_j^2\)，定义域要求每个 \(D_j^2u\) 存在于 \(L^2\)。
\optionalexerciseitem{plancherel-bandlimited-evaluation}设 \(f\in L^2(\R^d)\)、\(R>0\)，且 \(\widehat f\) 在 \(B(0,R)\) 外几乎处处为零。证明 \(f\) 有连续代表并估计 \(\|f\|_\infty\)。
\end{enumerate}
